% Placeholder frame used only until physical-device screenshots are available.
\definecolor{pendinggray}{HTML}{F3F4F6}
\newcommand{\pendingdemo}[2]{%
\begin{figure}[H]
\centering
\fcolorbox{black!35}{pendinggray}{%
\begin{minipage}[c][4.0cm][c]{0.82\textwidth}
\centering
\textbf{Pending live demonstration evidence}\\[6pt]
\small #1\\[8pt]
\footnotesize Replace this framed placeholder with the final device screenshot before submission.
\end{minipage}}
\caption{#2}
\end{figure}}

\section{Experimental Results and System Integration}

\subsection{Evaluation Protocol and Experimental Results}
Speaker verification is evaluated with 10,000 deterministic balanced trials on held-out VIVOS speakers: 5,000 genuine trials and 5,000 impostor trials. Given two embedding vectors $a$ and $b$, the runtime score is cosine similarity:
\[
    \mathrm{score}(a,b) = \frac{a \cdot b}{\lVert a \rVert_2\lVert b \rVert_2}.
\]
The report uses equal error rate (EER), the point where false reject and false accept rates are equal, and minimum detection cost function (minDCF). SID Top-1 accuracy is computed by averaging five enrollment embeddings for each test speaker, then assigning each remaining query recording to the highest-scoring profile.

\begin{table}[H]
\centering
\begin{tabularx}{\textwidth}{Xrrr}
\toprule
\textbf{System} & \textbf{EER} & \textbf{minDCF} & \textbf{SID Top-1} \\
\midrule
Untouched VoxCeleb ECAPA baseline & 4.40\% & 0.002978 & 100.00\% \\
One-epoch smoke model & 1.58\% & 0.001570 & 100.00\% \\
VIVOS fine-tuned ECAPA & \textbf{1.28\%} & \textbf{0.001628} & \textbf{100.00\%} \\
\bottomrule
\end{tabularx}
\caption{Held-out VIVOS evaluation results.}
\label{tab:results}
\end{table}

Fine-tuning lowers EER by 3.12 percentage points, or 70.9\% relative to the baseline. The experiment's EER operating threshold is 0.34559793992808. It is not automatically deployed: the current application threshold remains 0.80 and requires calibration on representative elderly-user recordings before it can be changed. The test-set result therefore demonstrates Vietnamese-domain performance, not a clinical or elderly-population validation.

\subsection{Speaker Enrollment, Verification, and Identification Integration}
Enrollment is available in Android settings. The user reads five prompts; the client captures each recording and posts it with its phrase index to the backend. The backend forwards the audio to the FastAPI enrollment route, receives an ECAPA embedding, upserts the phrase record, and updates the profile average. A query uses stored profiles only from the authenticated caller and directly linked caregiver/elder accounts. The backend checks that an AI-returned identity belongs to that trusted set before exposing \texttt{recognized\_user\_name} for personalization.

Figure~\ref{fig:protected-flow} describes the protected medication flow. The voice grant is random, valid for two minutes, bound to a user and intent, usable once, and stored as a SHA-256 hash. The medication endpoint consumes the grant transactionally; a request without a valid grant is rejected.

\begin{figure}[H]
\centering
\begin{tikzpicture}[node distance=0.8cm and 0.45cm]
\node[service] (user) {Elder in Android app};
\node[service, right=of user] (query) {Voice query\\protected intent detected};
\node[service, right=of query] (verify) {Second recording\\ECAPA verification};
\node[service, below=of verify] (grant) {One-use voice grant\\user + intent + expiry};
\node[service, left=of grant] (med) {Mark dose taken\\grant consumed};
\node[datastore, below=of med] (database) {PostgreSQL};
\draw[flow] (user) -- (query);
\draw[flow] (query) -- (verify);
\draw[flow] (verify) -- (grant);
\draw[flow] (grant) -- (med);
\draw[flow] (med) -- (database);
\end{tikzpicture}
\caption{Speaker-verification-gated medication action. A first query never authorizes the mutation.}
\label{fig:protected-flow}
\end{figure}

\subsection{Implementation Validation and Demonstration Evidence}
The completed engineering checks include AI-service linting and focused tests, Docker Compose configuration validation, a ready AI container with the trained checkpoint, backend unit tests for voice grants and general responses, and a successful Android debug build. These are implementation checks; they do not replace a physical-device demonstration. In particular, the current evidence does not claim successful Firebase authentication, microphone capture, Vietnamese TTS playback, real enrollment, SV, or SID on a device.

The following captioned placeholders reserve the exact visual evidence required for the final submission. They must be replaced with real screenshots or video frames before the report is submitted.

\pendingdemo{Android settings showing the five-phrase enrollment dialog and a completed $5/5$ state.}{Reserved evidence: speaker enrollment workflow.}
\pendingdemo{A recorded general request for time, date, or the medication schedule, with its visible response.}{Reserved evidence: unauthenticated general voice function.}
\pendingdemo{Successful medication verification followed by a second-speaker failure case.}{Reserved evidence: SV-protected medication action.}
\pendingdemo{Two enrolled and linked users receiving a response personalized with the recognized name.}{Reserved evidence: SID personalization.}

\subsection{Limitations and Conclusion}
Dearly provides a reportable end-to-end design: VIVOS fine-tuning improves held-out verification performance, ECAPA embeddings drive both SV and SID, and a server-side grant prevents a classified protected voice command from becoming an authorization decision. The remaining limitation is evidence collection rather than a claim of completed live validation. The final package must include real-device demonstration media, source code, the report, and a Drive-link text file for large dataset/model artifacts when needed. The ZIP filename for this team should use both student IDs: \texttt{23127118\_23127144.zip}.
